\section{练习题}
先独立完成，再看题后解答。原题单十五道第二题中，五道在真题速看展开，其余十道在本节提供解答与完整实现。先按原题单完成，再做算法专项补充。
\subsection{梦境巡查}
\textbf{第36次CCF-CSP认证，2024-12，第2题。}\csp{CSP202412B}{梦境巡查}。
\textbf{题目与输入输出。}输入n，接着输入$a_0,\ldots,a_n$及$b_1,\ldots,b_n$。从0号区域出发，走第i段消耗$a_i$，到区域i后得到$b_i$；每段出发前能量须够支付。当仅第i个补给失效时，求完成整条路线并返回0号区域所需的最少初始能量$w_i$，输出n个结果。$1\le n\le10^5$，所有消耗与补给在0到1000之间。
\textbf{评测限制。}曙梦OJ当前时限1秒，内存512MB。
\textbf{读题思路。}n次逐段模拟为$O(n^2)$，规模不允许。缺失$b_i$只影响从区域i出发及以后的检查；区域i之前不受影响。
\textbf{做题思路与解答。}设$d_j=\sum_{t=0}^j a_t-\sum_{t=1}^j b_t$，表示出发走第j段前，对初始能量的下界。缺失第i个补给后，$j<i$仍需$d_j$，$j\ge i$则需$d_j+b_i$。分别预处理d的前缀最大值P和后缀最大值S，得到$w_i=\max(P_{i-1},S_i+b_i)$。因为$d_0=a_0\ge0$，无需额外把负答案截成0。每个出发检查都满足当且仅当初始值不低于这些下界的最大值，因此这是最小可行值。
\begin{verbatim}
输入
3
5 5 5 5
0 100 0
输出
10 20 10
\end{verbatim}
\begin{verbatim}
输入
3
9 4 6 2
9 4 6
输出
15 10 9
\end{verbatim}
\textbf{复杂度与边界。}时间$O(n)$，空间$O(n)$。检查b本来就是0和巨额补给位于中间的情况。
\Code{45_dream.cpp}
\subsection{机器人宿管指南}
\textbf{第42次CCF-CSP认证，2026-05，第2题。}\csp{CSP202605B}{机器人宿管指南}。
\textbf{题目与输入输出。}输入初始苹果数n、损耗百分比k、天数m。每天先丢掉$\lceil tk/100\rceil$个坏苹果（t是前一天剩余量），再给每个机器人一个苹果。求最多能供养多少机器人，要求第m天吃完后剩余非负。$1\le n\le10^{18}$，$1\le k\le5$，$1\le m\le2000$；保证答案为不超过$10^9$的正整数。
\textbf{评测限制。}曙梦OJ当前时限0.5秒，内存512MB。
\textbf{读题思路。}顺序是先变质后食用，百分比向上取整。n超过int范围，即使long long也不宜直接算$n\times100$；不能用浮点数近似大整数比例。
\textbf{做题思路与解答。}人数x给定后按天模拟检查。剩余映射为$\lfloor t(100-k)/100\rfloor-x$，对t单调不减，所以人数越少，全程剩余越多，可行人数形成前缀。二分最后一个可行人数，范围$[0,\min(10^9,n/m)]$。损耗按$\lfloor t/100\rfloor k+\lceil (t\bmod100)k/100\rceil$计算，各项都在long long范围内。
\begin{verbatim}
输入
1000 1 5
输出
194
\end{verbatim}
\begin{verbatim}
输入
19 5 3
输出
5
\end{verbatim}
\textbf{复杂度与边界。}时间$O(m\log 10^9)$，空间$O(1)$。0作为可行哨兵，题面保证最终答案为正。第m天可以恰好剩0，不要求还能吃第m+1天。
\Code{46_dorm.cpp}
\subsection{字符串变换}
\textbf{第35次CCF-CSP认证，2024-09，第2题。}\csp{CSP202409B}{字符串变换}。
\textbf{题目与输入输出。}字符集合为大小写字母、数字和空格，共63种。输入以\#包围的初始字符串，再给n个形如\#xy\#的字符替换规则，未定义的字符保持不变；规则构成单射。随后输入m与m个k，每次独立询问初始串连续变换k次的结果，以\#包围输出。$n\le63,m\le1000,k\le10^9$，初始串长不超过100。
\textbf{评测限制。}曙梦OJ当前时限1秒，内存512MB。
\textbf{读题思路。}保留空格，不能用\%s读整行；每次查询从原始串开始。每轮替换是同时进行，不能按规则就地多次替换同一个位置。巨大k不能逐轮模拟。
\textbf{做题思路与解答。}预处理$jump[b][c]=f^{2^b}(c)$，初值为一次变换，递推$jump[b][c]=jump[b-1][jump[b-1][c]]$。把k拆成二进制，对每个字符依次应用对应的跳表即可。每个字符独立变化，函数复合保证跳表恰好等于所需次数。单射还允许按置换环长度取模；这里的倍增实现无需依赖单射，便于统一处理。
\begin{verbatim}
输入
#Hello World#
6
#HH#
#e #
# r#
#re#
#oa#
#ao#
3
1 2 3
输出
#H llarWaeld#
#HrlloeWo ld#
#Hella Warld#
\end{verbatim}
\textbf{复杂度与边界。}字符表用128项覆盖ASCII；时间$O(128\log K+m|s|\log K)$，空间$O(128\log K+|s|)$。输入的两个\#是包装符号，不参与变换。
\Code{47_transform.cpp}

\subsection{机器人复健指南}
\textbf{第38次CCF-CSP认证，2025-06，第2题。}\csp{CSP202506B}{机器人复健指南}。
\textbf{题目与输入输出。}输入n、k和起点x、y。在$n\times n$场地内每次按马步跳跃，即坐标变化为$(\pm1,\pm2)$或$(\pm2,\pm1)$，共八个方向，不得跳出场地。统计不超过k步可到达的不同格数，包含起点。$1\le n,k\le100$。输出格数。
\textbf{读题思路。}图中的八方向是马步，不是上下左右及四个对角相邻格。要求不超过k步，且一个格子多次到达只计一次。
\textbf{做题思路与解答。}BFS把每个格子作为状态，首次访问得到最短步数。到达第k层就不再扩展；每个首次入队的格子给答案加1，起点先计入。无权边BFS按距离递增发现状态，所以所有且仅有最短距离不超过k的格子被统计。也能用“步数、位置”的分层可达DP，再对各层取并集，但BFS不用重复访问位置，复杂度更低。
\begin{verbatim}
输入
4 1
1 1
输出
3
\end{verbatim}
\begin{verbatim}
输入
4 2
1 1
输出
8
\end{verbatim}
\textbf{复杂度与边界。}时间$O(n^2)$，空间$O(n^2)$，每格最多一次入队。n=1或2时可能完全没有合法马步；答案仍包含起点。
\Code{48_knight.cpp}

\subsection{回收站选址}
\textbf{第18次CCF-CSP认证，2019-12，第2题。}\csp{CSP201912B}{回收站选址}。
\textbf{题目与输入输出。}输入n个两两不同的垃圾点坐标。只有本身有垃圾且上下左右四邻都存在的点可以选址；得分为四个对角邻点中有垃圾的个数。依次输出得分0到4的选址数量，各一行。$n\le1000$，坐标绝对值不超过$10^9$。
\textbf{读题思路。}地图坐标巨大且可负，不能按坐标开完整网格。只有已有垃圾点能成为候选；四个正交邻居是资格条件，四个对角邻居只负责评分。
\textbf{做题思路与解答。}用set保存坐标pair。枚举每个原点，查询四个正交邻点，全部存在后统计四个对角邻点并增加相应频率。set的键同时包含横纵坐标，查询精确匹配，所以资格和得分都符合定义。n很小时也可用数组线性查找作为对照。
\begin{verbatim}
输入
7
1 2
2 1
0 0
1 1
1 0
2 0
0 1
输出
0
0
1
0
0
\end{verbatim}
\begin{verbatim}
输入
2
0 0
-100000 10
输出
0
0
0
0
0
\end{verbatim}
\textbf{复杂度与边界。}时间$O(n\log n)$，空间$O(n)$。输入坐标绝对值至多$10^9$，加减1仍在int范围内；n为1时输出五个0。
\Code{58_recycle.cpp}
\subsection{风险人群筛查}
\textbf{第20次CCF-CSP认证，2020-09，第2题。}\href{https://www.acwing.com/problem/content/3296/}{风险人群筛查（在线做题）}。
\textbf{题目与输入输出。}输入n、k、t及矩形左下xl、yd、右上xr、yu；再输入每人的t个记录坐标。含边界的矩形内出现过记录即“经过”；有至少k条连续记录都在内即“逗留”。输出经过人数和逗留人数，各一行。数据范围$n\le24$，$1\le k\le t\le1000$，坐标绝对值不超过$10^6$。
\textbf{读题思路。}只检查采样记录，不插值判断相邻记录间的连线。逗留强调连续；分散的k次进入不算。矩形边界也在区域内。
\textbf{做题思路与解答。}对每人维护曾进入标记、当前连续进入长度run、曾逗留标记。内点令run加1，外点清0；run达到k后保留逗留标记，即使后来离开也不取消。每人结束时各标记只加一次计数。
\begin{verbatim}
输入
1 3 8 0 0 10 10
-1 -1 0 0 0 0 -1 -1 0 0 -1 -1 0 0 0 0
输出
1
0
\end{verbatim}
\textbf{复杂度与边界。}时间$O(nt)$，额外空间$O(1)$。已经确定逗留后仍须读取剩余坐标，否则下个人输入错位。
\Code{59_risk.cpp}
\subsection{非零段划分}
\textbf{第23次CCF-CSP认证，2021-09，第2题。}\href{https://www.acwing.com/problem/content/4010/}{非零段划分（在线做题）}。
\textbf{题目与输入输出。}输入n项非负整数A，选择正整数p，把所有小于p的值变0，求剩余连续非零段数量的最大值。一个非零段必须不能再向左或向右扩展。$n\le5\times10^5,A_i\le10^4$。输出最大段数。
\textbf{读题思路。}阈值变化时段数不单调，不能二分“段数达到某值”。直接枚举每个p并扫描全数组约$5\times10^9$次，需利用值域。
\textbf{做题思路与解答。}每段恰好由一个开头贡献。设$A_0=0$，位置i成为开头当且仅当$A_{i-1}<p\le A_i$。因此当$A_{i-1}<A_i$时，给阈值区间$[A_{i-1}+1,A_i]$贡献1，用差分记录；遍历阈值恢复段数并取最大值。每个真实段只有一个开头，条件与开头定义等价，所以统计不重不漏。
\begin{verbatim}
输入
11
3 1 2 0 0 2 0 4 5 0 2
输出
5
\end{verbatim}
\begin{verbatim}
输入
14
5 1 20 10 10 10 10 15 10 20 1 5 10 15
输出
4
\end{verbatim}
\textbf{复杂度与边界。}时间$O(n+V)$，V=$10^4$，空间$O(V)$。数组全0时答案0；为$A_i+1$处预留容量。
\Code{60_nonzero.cpp}
\subsection{出行计划}
\textbf{第25次CCF-CSP认证，2022-03，第2题。}\href{https://www.acwing.com/problem/content/4458/}{出行计划（在线做题）}。
\textbf{题目与输入输出。}输入计划数n、查询数m、等待结果时长k；随后n行出行时刻$t_i$和有效时长$c_i$，最后m个检测时刻q。q检测的结果可从$q+k$到$q+k+c_i-1$进入对应场所。每查询输出符合条件的计划数。$n,m,k\le10^5$，$t_i,c_i,q\le2\times10^5$，均为正整数。
\textbf{读题思路。}有效时间包括刚获得结果的时刻，末端是有效时长减1。把“这次检测满足哪些计划”反过来，求“每项计划允许在哪些时刻检测”。
\textbf{做题思路与解答。}条件$q+k\le t_i\le q+k+c_i-1$等价于$t_i-c_i+1-k\le q\le t_i-k$。每项计划给这个整数区间加1，按q正数范围裁剪后做差分；一次前缀累加得到每个q的答案。等价变形保证区间内恰好计入满足要求的计划。
\begin{verbatim}
输入
6 2 10
5 24
10 24
11 24
34 24
35 24
35 48
1
2
输出
3
3
\end{verbatim}
\textbf{复杂度与边界。}时间$O(n+m+V)$，V=$2\times10^5$，空间$O(V)$。右端为负或裁剪后左端大于右端时跳过；r+1必须有空间。
\Code{61_travel.cpp}
\subsection{寻宝！大冒险！}
\textbf{第26次CCF-CSP认证，2022-06，第2题。}\csp{CSP202206B}{寻宝！大冒险！}。
\textbf{题目与输入输出。}地图是坐标0到L的方格点，共n棵树，输入n、L、S及n个树坐标，再输入$(S+1)\times(S+1)$的0/1藏宝图。图的第一维按S到0输入，第二维按0到S输入。藏宝图左下格$B[0][0]=1$。求所有与藏宝图完全匹配且不越地图边界的左下位置数。$n\le1000,L\le10^9,S\le50$，树坐标不重复。
\textbf{读题思路。}不能按L开地图，也不能枚举所有坐标。左下格保证有树，只需枚举n棵树为候选。匹配要求1与0都一致，不能只找到规定的树就忽略多余树。输入首行对应B[S]，不要按B[0]读。
\textbf{做题思路与解答。}先数藏宝图中的树数need。对每个候选，拒绝越界窗口；扫描全部地图树，若落入窗口却对应藏宝图0，立即拒绝，否则统计窗口树数。没有多余树且count=need时，一一对应保证所有藏宝图1也都有树，其余格为0，故整个窗口完全匹配。
\begin{verbatim}
输入
5 100 2
0 0
1 1
2 2
3 3
4 4
0 0 1
0 1 0
1 0 0
输出
3
\end{verbatim}
\begin{verbatim}
输入
5 4 2
0 0
1 1
2 2
3 3
4 4
0 0 0
0 1 0
1 0 0
输出
0
\end{verbatim}
\textbf{复杂度与边界。}时间$O(n^2+S^2)$，空间$O(n+S^2)$。代码严格按题面$B[S-i][j]$的输入下标存储；不额外转置。
\Code{62_treasure.cpp}
\subsection{何以包邮？}
\textbf{第27次CCF-CSP认证，2022-09，第2题。}\href{https://www.acwing.com/problem/content/4703/}{何以包邮？（在线做题）}。
\textbf{题目与输入输出。}输入n、包邮阈值x，随后n本书的正整数价格，每本只能选一次。求总价至少x的最小选购总价，保证全部价格和不小于x。$n\le30,a_i\le10^4$，故总价S不超过$3\times10^5$。输出最小总价。
\textbf{读题思路。}不是花费不超过x的最大价值，而是总价至少x且最小。$2^{30}$枚举太大；利用价格和上界做01可达性背包。
\textbf{做题思路与解答。}reachable[s]表示用已处理书本能否恰好买到s元。初值仅0为真，每本价格a加入时从S向a倒序更新$reachable[s]\mathrel{|}=reachable[s-a]$，保证当前书只使用一次。全部处理后从x向上找首个可达s，就是满足包邮的最小价格。
\begin{verbatim}
输入
4 100
20
90
60
60
输出
110
\end{verbatim}
\begin{verbatim}
输入
3 30
15
40
30
输出
30
\end{verbatim}
\textbf{复杂度与边界。}时间$O(nS)$，空间$O(S)$。重复价格仍代表不同书，不能去重；全部购买可行保证扫描能找到答案。
\Code{63_shipping.cpp}

\subsection{机器人饲养指南：完全背包复做}
\textbf{第37次，2025-03，第2题。}\csp{CSP202503B}{机器人饲养指南}。
题面、样例和完整程序见本篇第三节。遮住答案重新写出$f[v]$，并回答：如果允许剩苹果，初始化和答案位置会有什么变化？\textbf{解答：}当前题要求恰好用完，不能把所有状态初始化为0。若改为最多吃n个且价值非负，可用“不超过容量”的全0初值，也可保留恰好状态并取$\max_{0\le v\le n}f[v]$。两种状态解释要与实现一致。
\subsection{数字排序：STL比较规则复做}
\textbf{第4次，2015-03，第2题。}\csp{CSP201503B}{数字排序}。
\textbf{任务。}输入n个整数，统计每个不同值的出现次数，按次数从大到小输出“值 次数”；次数相同按值从小到大输出。\textbf{解答与关键代码。}值域小时先用桶计数，再把非零记录存入结构体数组排序。值域大时可用map计数后生成记录。比较函数只表达严格顺序，不能用大于等于。
\begin{lstlisting}
struct Record { int value, count; };
bool less_record(const Record &a, const Record &b) {
    if (a.count != b.count) return a.count > b.count;
    return a.value < b.value;
}
// std::sort(records, records+total, less_record);
\end{lstlisting}
排序时间$O(u\log u)$，u是不同值个数；还需加上计数过程的成本。与第一篇众数题区别在于，这里输出全部不同值，并有两层排序键。
\subsection{其他算法专项训练}
以下题按本篇算法顺序安排。
\begin{longtable}{P{4.8cm}P{3.2cm}P{6cm}}
\toprule 做题链接 & 训练点 & 解答要点与边界\\\midrule\endhead
\csp{CSP201409B}{画图（第2次）} & 二维差分 & 每块矩形标记四角，恢复后统计覆盖数大于0的单位格，重叠只计一次。\\
\csp{CSP202403B}{相似度计算（第33次）} & string、set & 先统一大小写，每篇各自去重，统计交集与并集；同篇重复出现不能反复计数。\\
\csp{CSP202305B}{矩阵运算（第30次）} & 结合律、复杂度 & 对$QK^TV$先算$K^TV$，避免构造$n\times n$中间矩阵；输出还需按题意乘对应行权重。\\
\prob{P2678} 跳石头 & 二分、贪心 & 固定距离后从左到右判断最少移走数量，终点不可移走。\\
\prob{P1706} 全排列 & DFS & 选未使用元素，递归后撤销；原题要求每数宽度5。\\
\prob{P1443} 马的遍历 & BFS & 八个马步方向，首次入队记录最短距离，不可达输出-1。\\
\prob{B3637} 最长上升子序列 & 序列DP & 严格上升用$a_j<a_i$；大范围时对tails使用lower\_bound。\\
\prob{P1115} 最大子段和 & 递推DP & $f_i=\max(a_i,f_{i-1}+a_i)$；全负时不能把答案设成0。\\
\prob{P1048} 采药 & 01背包 & 每件一次，容量倒序；容量含义为最多使用。\\
\prob{P1616} 疯狂的采药 & 完全背包 & 每件不限次数，容量正序；价值总和与时限重新估算。\\
\prob{P1757} 通天之分组背包 & 分组背包 & 每组至多一件，不能让同组候选相互接续。\\
\bottomrule\end{longtable}

\subsection{算法补充真题：稀疏存储与基础数学}
以下两题来自历年考点补充，均在算法讲解旁给出做题入口。
\subsection{稀疏向量}
\textbf{第19次CCF-CSP认证，2020-06，第2题。}\csp{CSP202006B}{稀疏向量}。
\textbf{题目与输入输出。}长度为n的两个向量，只有a、b项非零。每个向量的非零项按下标严格递增输入，每项为下标、数值；求内积$\sum_{i=1}^n u_iv_i$。首行n、a、b，随后a行u和b行v。$n\le10^9$，$0<a,b<n$，$a,b\le5\times10^5$，非零数值绝对值不超过100。输出内积。
\textbf{评测限制。}曙梦OJ当前时限2秒，内存512MB。
\textbf{读题思路。}n很大，但非零项数较小，按n开数组会浪费数GiB。只有相同下标的两项才产生非零乘积，排序信息已经给出。最多约$5\times10^9$的绝对值累加，需要long long。
\textbf{做题思路与解答。}将每组存成下标和值的结构体数组。下标较小的那项在另一组当前及后续项中不可能匹配，可以丢弃；下标相同则乘积累加并同时前进。每个共同非零位置只遇到一次，其他位置贡献为0，所以扫描值等于原内积。
\begin{verbatim}
输入
10 3 4
4 5
7 -3
10 1
1 10
4 20
5 30
7 40
输出
-20
\end{verbatim}
\textbf{复杂度与边界。}时间$O(a+b)$，空间$O(a+b)$，与完整向量长度n无关。若改成映射也能做，但已有有序性时双指针更简单。
\Code{41_sparse.cpp}
\subsection{因子化简}
\textbf{第32次CCF-CSP认证，2023-12，第2题。}\csp{CSP202312B}{因子化简}。
\textbf{题目与输入输出。}对整数n进行质因数分解$n=\prod p_i^{e_i}$，删除所有指数$e_i<k$的整个因子$p_i^{e_i}$，保留其余因子，输出结果；全部删去时为1。输入q，随后q行n、k；每个结果一行。$1\le q\le10$，$1<n\le10^{10}$，$1<k\le10$。
\textbf{评测限制。}曙梦OJ当前时限1秒，内存512MB。
\textbf{读题思路。}“删除指数不足的整个质因子幂”不是把每个指数减少k，也不是只保留$p_i^k$。例如$360=2^3\cdot3^2\cdot5$，k=3时答案为8。n超过int上界，读入必须用long long。
\textbf{做题思路与解答。}从p=2开始试除，每找到一个质因子就反复整除，记录完整指数以及对应幂。指数达标才把该幂乘入答案。条件$p\le n/p$避免$p^2$溢出，也利用不断缩小的n减少试除。结束后残留大于1的数是指数为1的质因子；由于本题k至少为2，它不保留。所有质因子都被处理且筛选严格按定义，得到所求乘积。
\begin{verbatim}
输入
3
2155895064 3
2 2
10000000000 10
输出
2238728
1
10000000000
\end{verbatim}
\textbf{复杂度与边界。}最坏时间$O(q\sqrt{\max n})$，空间$O(1)$。答案初值为1，覆盖所有因子被删除的情况。
\Code{43_factor.cpp}
